\begin{trapbox}
	\begin{description}[style=nextline, leftmargin=2.8em, labelsep=0.5em, itemsep=0.45em]
		\item[\textbf{\textcolor{CTrap}{易错点一：两点放在同侧}}]
		      若取 $A(m,a)$、$B(n,b)$，两点都在 $x$ 轴上方。
		      线段 $AB$ 不经过 $x$ 轴，不能让
		      $P(x,0)$ 在线段 $AB$ 上，因此不能直接把最小值写成 $|AB|$。

		      \textbf{改法：} 把其中一点关于 $x$ 轴对称。
		      例如取 $A(m,-a)$、$B(n,b)$，
		      对应的 $|PA|$、$|PB|$ 不变，但 $A$、$B$ 分居异侧。

		\item[\textbf{\textcolor{CTrap}{易错点二：只写“共线”}}]
		      三角形不等式的等号条件是 $P$ 在线段 $AB$ 上。
		      如果 $P$ 位于线段的延长线上，即使三点共线，
		      仍有 $|PA|+|PB|>|AB|$。

		      \textbf{本题为何能取等：} $a,b>0$，
		      所以线段 $AB$ 必穿过 $x$ 轴；令 $P$ 为交点即可。

		\item[\textbf{\textcolor{CTrap}{易错点三：纵坐标差算错}}]
		      $A$、$B$ 的纵坐标分别为 $-a$、$b$，
		      所以纵向距离是
		      \[
			      |b-(-a)|=a+b,
		      \]
		      不是 $|a-b|$。

		\item[\textbf{\textcolor{CTrap}{易错点四：忽略定义域限制}}]
		      公式给出的取等位置是
		      \[
			      x_0=\frac{bm+an}{a+b}.
		      \]
		      当 $x\in\mathbb{R}$ 时一定能够取到；
		      若题目另有限制，例如 $x\geq c$，
		      应先检查 $x_0$ 是否属于允许范围。
	\end{description}
\end{trapbox}